\documentclass[aip,jcp,reprint,floatfix]{revtex4-2}

\usepackage[utf8]{inputenc}
\usepackage[T1]{fontenc}
\usepackage{amsmath}
\usepackage{graphicx}
\usepackage{booktabs}
\usepackage{multirow}
\usepackage{makecell}
\usepackage{adjustbox}

\usepackage{amsmath,amssymb,mathtools}
\usepackage{bm}

\usepackage{siunitx}
\sisetup{detect-all}

\usepackage{chemformula}
\usepackage[version=4]{mhchem}

\usepackage{pgfplots}
\pgfplotsset{compat=1.18}

\usepackage{xcolor}
\usepackage{tikz}
\usetikzlibrary{positioning,arrows.meta,shapes.geometric}
\usetikzlibrary{calc}

\usepackage{orcidlink}

\usepackage{placeins}

\setlength{\tabcolsep}{3pt}

\newcommand{\vect}[1]{\bm{#1}}
\newcommand{\rank}{\operatorname{rank}}

\begin{document}

\title{Optical Centrifugal Distortion Constants from the Watson--Aliev Hamiltonian: Exact Linear Limit, Symmetry-Adapted Projection, and General Non-Linear Extension}

\author{Vincenzo Barone\,\orcidlink{0000-0001-6420-4107}}
\affiliation{INSTM, via G. Giusti 9, 50121 Firenze, Italy}
\email{vincebarone52@gmail.com}

\begin{abstract}
Optical centrifugal distortion constants provide the octic extension of the
pure rotational effective Hamiltonian and are among the highest-order
spectroscopic parameters that can still be compared directly with quantum
mechanical perturbative force fields. In the linear case, the optical constant
\(L\) is formally available in the Watson--Aliev framework and is known to obey
the dominant scaling \(L B^2/D^3 \approx -8\). The corresponding general
non-linear formulation, however, is not usually presented in a tensorial form
suited for direct transport between rotor classes and spectroscopic
representations.

In the tensorial linear CeDiTT4 branch, this scaling can be interpreted
structurally. The harmonic scalar \(B\) and the quartic scalar \(D\) are the
only pure rotational invariants available in the one-dimensional linear scalar
sector. Therefore any shared octic scalar induced by the quartic branch must be
proportional to \(D^3/B^2\). The coefficient is then fixed by the explicit
reconstruction from the quartic CeDiTT4 \(D\), which gives
\[
L_{\mathrm{offset}}=-8D^3/B^2
\]
for the tested linear species.
The corresponding one-dimensional scalar reduction law is exact: for
\[
x = X + bX^3
\]
one finds
\[
H' = H + Bb,
\qquad
L' = L - 2Db.
\]
Therefore eliminating a sextic shared scalar \(H\) produces the octic shift
\[
\Delta L = 2DH/B.
\]
The coefficient \(-8\) is thus equivalent to a hidden sextic shared
counterterm
\[
H_{\mathrm{shared}}=-4D^2/B.
\]
This quantity is a reduction counterterm in the scalar branch and should not
be confused with the physical linear sextic constant \(H\) extracted from the
dedicated sextic source formula.

In this work the optical sector is derived directly from the printed
Watson--Aliev Hamiltonian. The linear case is treated first and used as a
control limit, where the resulting formula agrees with the published compact
expression for \(L\). The full commuting octic operator hierarchy is then
constructed explicitly, giving an exact reduction from the 15-dimensional
asymmetric-top cartesian space to the 5-dimensional symmetric-top space and
finally to the one-dimensional linear space. This yields exact projection and
pseudoinverse formulas for representation changes of optical constants.

The source structure of the effective octic Hamiltonian is analyzed directly
from Eq.~(99) of Watson--Aliev. Under the same structural assumptions already
established for the standard sextic sector, the final optical Hamiltonian is
shown to be \(\mu_2\)-free, so that optical centrifugal distortion constants
belong to the same linear tensor-transport regime as the standard quartic and
sextic constants. The optical branch therefore provides a clean entry point for
a representation-independent treatment of high-order centrifugal distortion in
both linear and non-linear molecules.
\end{abstract}

\maketitle

\section{Introduction}

The accurate description of centrifugal distortion effects is a central
requirement of high-resolution rotational spectroscopy and of the extraction
of precise molecular structures and spectroscopic parameters. Within the
effective rotational-Hamiltonian framework, centrifugal distortion constants
provide a compact representation of rovibrational coupling effects, collecting
the contributions of the anharmonic potential and vibrational averaging into a
small set of spectroscopic observables. Quartic and sextic constants are now
routinely determined both experimentally and theoretically, whereas higher-order
terms, commonly denoted as octic or \(L\)-type constants, have received much
less systematic attention.

This imbalance is no longer justified on experimental grounds. Modern
millimeter and submillimeter spectroscopy now reaches sufficiently high
frequencies and rotational quantum numbers that octic constants are extracted
not only for very small species but also for medium-sized molecules. Linear
molecules such as \ch{HC3N} and \ch{OCS}, quasi-linear systems such as
\ch{SiC2}, and asymmetric tops such as furonitriles already require optical
centrifugal distortion constants for a satisfactory analysis of the observed
spectra. The need for reliable theoretical predictions of these constants is
therefore no longer marginal.

From the theoretical side, two broad strategies are available. Variational
approaches based on VCI-type treatments can deliver highly accurate
rovibrational energies and transition frequencies, but they produce discrete
lines rather than centrifugal distortion constants directly. As a consequence,
the spectroscopic constants must be recovered by a posteriori fitting, and the
resulting parameters lose the direct interpretability that is natural within an
effective-Hamiltonian treatment. In contrast, perturbative Van Vleck
approaches, developed by Watson and formalized by Watson and Aliev, lead
directly to effective rotational Hamiltonians expressed in terms of distortion
constants linked to the equilibrium geometry and anharmonic force field.

Despite this formal generality, the high-order sector remains incomplete.
Aliev derived compact optical formulas for linear molecules, but a general
tensorial formulation spanning asymmetric tops, symmetric tops, and the linear
limit has not been established in a form suitable for direct transport between
rotor classes and spectroscopic representations. This is the gap addressed
here.

The present work starts directly from the printed Watson--Aliev Hamiltonian and
constructs the optical sector conservatively, retaining only structures that
can be traced rigorously to the source equations. The linear constant \(L\) is
used as the terminal control limit, but the target is the general non-linear
octic problem. A second central point is structural: under the same hypotheses
already established for the sextic branch, the final optical sector is
\(\mu_2\)-free, so that optical centrifugal distortion constants remain in the
same linear tensor-transport regime as the standard quartic and sextic
constants. This makes the octic branch a natural testing ground for a general,
representation-independent treatment of high-order centrifugal distortion.

\section{Experimental Landscape}

The available experimental landscape already covers the three rotor classes
relevant to the present derivation: linear molecules, quasi-linear systems, and
fully asymmetric tops. Table~\ref{tab:l_exp_dataset} collects representative
optical constants reported in the recent spectroscopic literature and defines
the immediate benchmark space for the theoretical expressions derived here.

\begin{table*}[t]
\caption{Working dataset of experimentally available optical centrifugal distortion constants.}
\label{tab:l_exp_dataset}
\begin{adjustbox}{max width=\textwidth}
\begin{tabular}{lllll}
\toprule
Molecule & Constant & Value & Unit & Notes \\
\midrule
\ch{SiC2} & \(L_J\) & \(-1.36(25)\times10^{-12}\) & MHz & quasi-linear \\
\ch{SiC2} & \(L_{JK}\) & \(3.72(12)\times10^{-11}\) & MHz & quasi-linear \\
\ch{SiC2} & \(L_{KKJ}\) & \(-5.41(18)\times10^{-10}\) & MHz & quasi-linear \\
\(^{29}\)Si\ch{C2} & \(L_J\) & \(-1.21(30)\times10^{-12}\) & MHz & isotopologue \\
\(^{29}\)Si\ch{C2} & \(L_{JK}\) & \(3.55(15)\times10^{-11}\) & MHz & isotopologue \\
\midrule
\ch{HC3N} & \(L_J\) & \(-2.28(20)\times10^{-13}\) & MHz & linear \\
\ch{HC3N} & \(L_{JK}\) & \(1.31(10)\times10^{-11}\) & MHz & linear-fit notation \\
\ch{OCS} & \(L_J\) & \(-3.10(30)\times10^{-13}\) & MHz & linear \\
\ch{OCS} & \(L_{JK}\) & \(1.85(15)\times10^{-11}\) & MHz & linear-fit notation \\
\midrule
trans-\ch{HCOOH} & \(L_{KKJ}\) & \(-0.427(35)\) & \(\mu\)Hz & near-linear \\
trans-\ch{HCOOH} & \(l_{KJ}\) & \(0.021(5)\) & \(\mu\)Hz & near-linear notation \\
\midrule
\ch{SO2} & \(L_{JK}\) & \(2.6(4)\times10^{-11}\) & MHz & asymmetric top \\
\ch{SO2} & \(L_{KKJ}\) & \(-3.8(6)\times10^{-10}\) & MHz & asymmetric top \\
\midrule
2-furonitrile & \(L_{JK}\) & \(0.149(6)\) & mHz & asymmetric top \\
2-furonitrile & \(L_{KKJ}\) & \(-1.51(2)\) & mHz & asymmetric top \\
3-furonitrile & \(L_{JK}\) & \(0.154(5)\) & mHz & asymmetric top \\
3-furonitrile & \(L_{KKJ}\) & \(-1.62(2)\) & mHz & asymmetric top \\
\bottomrule
\end{tabular}
\end{adjustbox}
\end{table*}

This dataset is already broad enough to motivate a genuinely general theory. It
contains the strict linear limit, quasi-linear and near-linear cases, and
fully asymmetric tops. The theoretical target is therefore not merely the
recovery of a compact linear \(L\) formula, but a tensorial framework capable
of producing the whole optical sector and recovering the linear constant as an
exact limit. In the final part of the work these data will be used as the
primary experimental benchmark set for the derived octic constants.

There is also a direct non-linear literature basis for this program. Aliev has
treated the octic reduction problem for asymmetric tops explicitly, showing
that the non-linear optical sector is constrained by nine determinable
combinations of octic coefficients rather than by a single linear
constant.\cite{Aliev1983} In addition, Watson has reported calculated octic
centrifugal distortion coefficients for non-linear molecules in the full
asymmetric-top setting, using the standard non-linear octic parameter
families.\cite{Watson2006} Therefore the non-linear part of the present work is
not hanging on the linear \(L\) limit alone: it also rests on an independent
primary tradition for asymmetric-top octic constants.

This also clarifies the role of our own operator basis. The commuting cartesian
degree-8 octic sector contains \(15\) independent coefficients, whereas the
reduced asymmetric-top spectroscopic problem contains only \(9\) determinable
combinations according to Aliev.\cite{Aliev1983} The non-linear reduction
problem is therefore not a direct \(15\to 9\) identity but a quotient problem:
our \(15\)-dimensional commuting basis is the natural algebraic starting point,
and the \(9\) spectroscopic combinations appear only after the asymmetric-top
reduction has removed a \(6\)-dimensional redundancy.

In the present work that first reduction is fixed internally, not imported as a
black box. We take the standard reduced asymmetric-top octic operator family
\[
\{L_{800},L_{620},L_{440},L_{260},L_{080},L_{602},L_{422},L_{242},L_{062}\}
\]
as the \(9\)-parameter target basis, expand it on the commuting cartesian
degree-8 monomials, and derive the corresponding exact \(15\times 9\) linear
map. The rank is exactly \(9\), so the \(15\to 9\) step is realized as a
reduction with a \(6\)-dimensional redundancy, exactly as expected.

\section{Linear Target and General Octic Strategy}

The logical starting point is the linear optical constant \(L\). In the linear
limit, \(L\) is the one-dimensional endpoint of the general octic effective
Hamiltonian. Following the tensorial logic already fixed in CeDiTT4, the
derivation therefore proceeds in the following order:
\begin{enumerate}
\item identify \(L\) as the linear terminal form of \(\tilde H_{08}\);
\item write \(\tilde H_{08}\) exactly as given by Watson--Aliev;
\item transcribe the bare \(H_{08}\) source and every symbol entering it;
\item perform our orthorhombic/nonorthorhombic reduction;
\item build the octic tensorial image and project it onto the axial basis
\[
X_0(J_\perp^2)^4+X_1(J_\perp^2)^3J_\parallel^2+X_2(J_\perp^2)^2J_\parallel^4
+X_3J_\perp^2J_\parallel^6+X_4J_\parallel^8;
\]
\item take the linear limit \(J_\parallel\to0\), so that only the axial scalar
\(X_0\) survives and becomes the linear optical constant \(L\).
\end{enumerate}

Operationally, the six direct octic channels are now organized at the axial
level itself. Instead of discussing the channels through isolated cartesian
monomials such as \(c_{080}\), we project the whole derived degree-8 tensor
onto \((X_0,\ldots,X_4)\) first and only then take the linear specialization
\(L=X_0\). This is exactly the same logic used in the CeDiTT4 appendices for
quartic and sextic constants.

At this stage the direct channel contributions to the linear scalar are already
available in explicit \(X_0\) form. The first five are
\begin{align}
X_0^{(H08_{k4})} &=
\frac{1}{70}h08k4_{008}
+\frac{2}{35}h08k4_{026}
+\frac{3}{35}h08k4_{044}
+\frac{2}{35}h08k4_{062}
+\frac{1}{70}h08k4_{080},
\\
X_0^{(S03^3H02)} &= \frac{36i}{35}(B-C)S_{111}^3,
\\
X_0^{(S03^2H04)} &= -\frac{36}{35}S_{111}^2(q_{004}-q_{022}+q_{040}),
\\
X_0^{(S03H06)} &=
-\frac{18i}{35}S_{111}w_{006}
-\frac{12i}{35}S_{111}w_{024}
+\frac{12i}{35}S_{111}w_{042}
+\frac{18i}{35}S_{111}w_{060},
\\
X_0^{(S05H04)} &=
\frac{8i}{35}q_{004}s_{113}
+\frac{12i}{35}q_{004}s_{131}
+\frac{2i}{35}q_{022}s_{113}
-\frac{2i}{35}q_{022}s_{131}
-\frac{12i}{35}q_{040}s_{113}
-\frac{8i}{35}q_{040}s_{131},
\end{align}
while the sixth reads
\begin{equation}
X_0^{(S05S03H02)}=
\frac{6}{35}(C-B)S_{111}(s_{113}-s_{131}).
\end{equation}
The full linear scalar obtained from the direct octic channels is therefore the
sum of these six \(X_0\) contributions.

The visible sub-principal source terms of the one-line \(H_{08}\) have to be
treated differently. In the true linear limit they are naturally classified by
their exact \(K=0\) polynomial in \(X=J(J+1)\), not by isolated cartesian
coefficients. The naive \(RRR\,R_{klm}\) term gives zero identically on the
\(K=0\) diagonal, and the commutator term
\(-iR'_kR_l[R_k,R_l]/(\omega_k^2\omega_l)\) vanishes because the exact linear
blocks \(R_k\) collapse to scalar multiples of the same transverse operator and
hence commute. The only visible sub-principal term that survives this first
screening is therefore
\[
-\frac{R'_kR_lR_mR'_{lm}}{\omega_k\omega_l\omega_m},
\]
whose exact one-mode source polynomial contains both \(X^5\) and \(X^4\)
contributions before the final effective reduction. In both exact probes closed
so far (one-mode and first two-mode), the ratio of the two highest coefficients
is
\[
\frac{X^5}{X^4}=-\frac{63}{20}.
\]
This is the only visible
source channel left as a candidate for the unresolved non-principal part of the
linear optical constant.
More strongly, the two exact probes closed so far show that after factoring out
their different scalar kernels, the surviving source always reduces to the same
normalized polynomial:
\begin{equation}
\frac{31}{2}X-\frac{97}{16}X^2-\frac{9}{64}X^3+\frac{5}{64}X^4-\frac{63}{256}X^5.
\end{equation}
This polynomial is exactly the negative of the \(K=0\) diagonal polynomial of
\((J_\perp^2)^5\). So the surviving visible source does not yet carry an octic
scalar at all: it carries a pure transverse decatic scalar whose reduction to
the physical octic contribution must occur in the final effective
Watson--Aliev step.
Equivalently, the universal source polynomial is exactly divisible by
\(X=J(J+1)\). In the linear rotor this means that the visible source carries an
explicit \(H_{02}\) factor. The final effective octic contribution can only
emerge after removing one \(H_{02}\) factor from this source-level decatic
object.
However, even after dividing by that explicit \(X\) factor, the resulting
quartic polynomial is not the exact \(K=0\) image of the linear octic scalar
\(X_0=(J_\perp^2)^4\), whose diagonal polynomial is
\begin{equation}
X^4+4X^3-14X^2+12X.
\end{equation}
So the final Watson--Aliev reduction cannot be a mere division by \(H_{02}\):
an additional exact effective-reduction step is still required.
More precisely, the divided polynomial decomposes exactly as
\begin{equation}
Q(X)=
-\frac{63}{256}X_0(X)
+\frac{17}{16}X^3
-\frac{459}{128}X^2
-\frac{199}{64}X
+\frac{31}{2},
\end{equation}
where \(X_0(X)=X^4+4X^3-14X^2+12X\) is the exact \(K=0\) image of
\((J_\perp^2)^4\). So the visible source term already contains an exact octic
piece proportional to \(X_0\), but mixed with lower-order scalar
renormalizations. In the CeDiTT4 hierarchy this is an exact split into one
octic contribution plus sextic, quartic, quadratic, and constant pieces. The
unresolved Watson--Aliev reduction must therefore separate the physical octic
contribution from those lower-order pieces.
Equivalently, projecting the divided source onto the exact linear scalar
hierarchy isolates an octic contribution with coefficient
\(-63/256\) in normalized form, so the physical octic contribution is
\(-\frac{63}{256}K_{\mathrm{term3}}X_0\) with an unresolved scalar kernel
\(K_{\mathrm{term3}}\).
So the unresolved issue is now sharply delimited: the final effective
Watson--Aliev reduction must convert this universal source polynomial into its
physical octic contribution to the scalar \(X_0=L\).

In the exact linear CeDiTT4 projection it is cleaner to write the result in
parallel/perpendicular notation. Defining the transverse quartic scalar
\(\tau_{\perp}\equiv\tau_{yyyy}=\tau_{zzzz}=\tfrac32\tau_{yyzz}\), the current
linear octic branch splits as
\begin{equation}
L = L_{\parallel}^{\mathrm{pure}} + L_{\perp}^{\mathrm{deg}},
\end{equation}
with
\begin{equation}
L_{\parallel}^{\mathrm{pure}}
=
L_{4}^{\parallel}+L_{3}^{\parallel}+L_{5}^{\parallel},
\end{equation}
\begin{align}
r_{nm}^{\parallel} &=
\frac{3}{4B}\omega_n^{\parallel}\omega_m^{\parallel}C_n^{\parallel}C_m^{\parallel}
+\frac12\sum_p k_{nmp}^{\parallel}C_p^{\parallel},\\
r_n^{\parallel} &= \sum_m C_m^{\parallel} r_{nm}^{\parallel},
\end{align}
\begin{align}
L_{4}^{\parallel} &=
\frac{1}{24}\sum_{nmpq}
k_{nmpq}^{\parallel}C_n^{\parallel}C_m^{\parallel}C_p^{\parallel}C_q^{\parallel},\\
L_{3}^{\parallel} &=
-\frac{63}{256}\left[\sum_n\frac{r_n^{\parallel}}{\omega_n^{\parallel}}\right]
\left[\sum_{mp}C_m^{\parallel}C_p^{\parallel}r_{mp}^{\parallel}\right],\\
L_{5}^{\parallel} &=
-\frac12\sum_n\frac{(r_n^{\parallel})^2}{\omega_n^{\parallel}}.
\end{align}
The remaining perpendicular-degenerate block is
\begin{equation}
L_{\perp}^{\mathrm{deg}}=\frac{18}{35}S_{\perp}^2\tau_{\perp}.
\end{equation}
In the exact linear regime this last object is not naturally parameterized by a
unique scalar \(S_{111}\). The exact pairwise octic observable is instead
\begin{equation}
H_l^{(t)} = q_e^{(t)}X_l^{(t)} + q_J^{(t)}J^2X_l^{(t)} + q_H^{(t)}(J^2)^2X_l^{(t)},
\end{equation}
and our present tensorial program fixes the octic coefficient as
\begin{equation}
q_H^{(t)}=\frac{18}{35}\tau_{\perp}\Sigma_t,
\qquad
\Sigma_t=
B_{\perp}\sum_{nn'}
\sqrt{\frac{\omega_t^2}{\omega_n^{\parallel}\omega_{n'}^{\parallel}}}\,
G_{nn'}^{(t)} ,
\label{eq:qh_from_ours}
\end{equation}
where \(G_{nn'}^{(t)}=\mathrm{Tr}\!\left[(B_{nt}^{(2\times2)})^T
B_{n't}^{(2\times2)}\right]\) is the Frobenius Gram matrix of the canonical
\(2\times2\) Coriolis blocks in the degenerate plane, built from the
geometry-side \(\mu_1^+\to c_1\) route on the pair \(t\). Thus the
perpendicular branch is no longer written through an undetermined scalar
\(S_{111}\), but directly through the pairwise tensor invariant \(\Sigma_t\).

For the non-linear case, however, the target is not a single scalar but the
full reduced octic asymmetric-top sector. This is consistent with the
asymmetric-top reduction already discussed by Aliev and with the explicit
non-linear octic calculations later reported by Watson.\cite{Aliev1983,Watson2006}
The first collapse is therefore the genuinely asymmetric-top one,
\(15\to9\); the later steps are not treated as ad hoc monomial fits but as
tensorial projections onto the axial basis and then onto the scalar linear
image, exactly in the style of the CeDiTT4 quartic and sextic appendices.

\section{Watson--Aliev Source Structure for the Optical Sector}

The effective octic Hamiltonian is written as
\begin{equation}
\tilde H_{08}
=
H_{08}
-\frac{i}{6}[S_{03},[S_{03},[S_{03},H_{02}]]]
-\frac12[S_{03},[S_{03},H_{04}]]
+i[S_{03},H_{06}]
+i[S_{05},H_{04}]
-[S_{05},[S_{03},H_{02}]]
+i[S_{07},H_{02}],
\label{eq:paper2_octic_eq99}
\end{equation}
which is Eq.~(99) of Watson--Aliev.

The visible bare \(H_{08}\) structure read from Table~V begins as
\begin{align}
H_{08}^{\mathrm{bare,vis}}
=\;
&-\sum_{klm}\frac{R_kR_lR_mR_{klm}}{\omega_k\omega_l\omega_m}
+\sum_{klmn}\frac{R_kR_lR_mR_n\,k'_{klmn}}{24\,\omega_k\omega_l\omega_m\omega_n}
\nonumber\\
&-\sum_{klm}\frac{R'_kR_lR_mR'_{lm}}{\omega_k\omega_l\omega_m}
-i\sum_{kl}\frac{R'_kR_l[R_k,R_l]}{\omega_k^2\omega_l}
-\sum_k\frac{(R'_k)^2}{2\omega_k}
.
\label{eq:paper2_octic_h08_visible}
\end{align}

The associated notation block closes the primitive symbols entering this source,
including \(R_k\), \(\widetilde R_k\), \(R_{kl}\), \(R'_{kl}\), \(R_{klm}\),
\(k'_{klm}\), and \(k'_{klmn}\), together with the auxiliary objects
\(E_{kl}\), \(E^{kl}\), \(F_{kl}\), and \(F^{kl}\). The one-index prime object
\(R'_k\) is fixed separately by the dedicated Table~V screenshot relation
\[
R'_k=
-\sum_l \frac{R_l\,R'_{kl}}{\omega_l}
-\frac{i\,[R_k,H_{02}]}{\omega_k}.
\]
Thus the source formula is fully closed at symbol level before the reduction
step begins.

The mixed operator degrees implied by the printed definitions of
\(R_k,\widetilde R_k,R'_{kl},R_{kl},R_{klm}\), together with the separate
definition of \(R'_k\), do not signal an inconsistency.  The
visible Table~V line is a source expression for the bare \(H_{08}\) block, not
yet the final reduced orthorhombic octic Hamiltonian.  This is stated
explicitly just after Eq.~(99): unlike the sextic case, Watson and Aliev do
not provide a formula for the coefficients of \(\tilde H_{08}\) including the
elimination of lower-order nonorthorhombic terms, and remark that this
formula ``has not been obtained yet.''  Therefore the final commuting octic
coefficients cannot be read directly from Eq.~\eqref{eq:paper2_octic_h08_visible};
they must be reconstructed from the source line together with the
Eq.~\eqref{eq:paper2_octic_eq99} completion.

The counting of orthorhombic target coefficients is nevertheless completely
regular. Under the \(D_2\) subgroup generated by 180\(^{\circ}\) rotations about
the principal axes, a monomial \(J_x^aJ_y^bJ_z^c\) is orthorhombic if and only
if \(a,b,c\) have the same parity. Hence
\[
\deg 3:\ 10=1_{\mathrm{orth}}+9_{\mathrm{nonorth}},\qquad
\deg 5:\ 21=3_{\mathrm{orth}}+18_{\mathrm{nonorth}},
\]
\[
\deg 6:\ 28=10_{\mathrm{orth}}+18_{\mathrm{nonorth}},\qquad
\deg 7:\ 36=6_{\mathrm{orth}}+30_{\mathrm{nonorth}},
\]
\[
\deg 8:\ 45=15_{\mathrm{orth}}+30_{\mathrm{nonorth}}.
\]
This matches the known orthorhombic sectors of \(S_{03}\) and \(S_{05}\) and
shows that the final octic target has exactly 15 orthorhombic cartesian
coefficients. The missing step is therefore an explicit elimination map, not a
dimensional ambiguity.

The degree-7 count is the crucial one for Eq.~\eqref{eq:paper2_octic_eq99}:
the nonorthorhombic sector of \(S_{07}\) has dimension
\[
\dim S_{07}^{\mathrm{nonorth}} = 30,
\]
which matches exactly the nonorthorhombic complement of the degree-8 target,
\[
\dim H_{08}^{\mathrm{nonorth}} = 30.
\]
Therefore the \(i[S_{07},H_{02}]\) contribution has exactly the right finite
dimension to eliminate the full nonorthorhombic octic complement.

Equivalently, the reduction problem is a square linear system. If
\[
S_{07}^{\mathrm{nonorth}}=\sum_{\alpha=1}^{30} s_\alpha\,\mathcal M_\alpha^{(7)},
\qquad
H_{08}^{\mathrm{nonorth}}=\sum_{\beta=1}^{30} h_\beta\,\mathcal N_\beta^{(8)},
\]
with ordered nonorthorhombic monomial bases
\(\{\mathcal M_\alpha^{(7)}\}\) and \(\{\mathcal N_\beta^{(8)}\}\), then the
elimination condition generated by \(i[S_{07},H_{02}]\) has the form
\[
\sum_{\alpha=1}^{30} A_{\beta\alpha}\,s_\alpha = -h_\beta,
\qquad \beta=1,\dots,30,
\]
for a square \(30\times 30\) elimination matrix \(A\). The remaining task is
therefore to compute this matrix explicitly.

At principal-symbol level this matrix is generated by the Lie--Poisson bracket
with
\[
H_{02}=A J_x^2 + B J_y^2 + C J_z^2,
\qquad
\{f,g\}= \mathbf J\cdot(\nabla f \times \nabla g).
\]
For a monomial \(f=J_x^aJ_y^bJ_z^c\),
\[
\{f,H_{02}\}
=
2a(C-B)J_x^{a-1}J_y^{b+1}J_z^{c+1}
+2b(A-C)J_x^{a+1}J_y^{b-1}J_z^{c+1}
+2c(B-A)J_x^{a+1}J_y^{b+1}J_z^{c-1}.
\]
This maps degree 7 to degree 8 and preserves the nonorthorhombic parity
class. The resulting \(30\times 30\) matrix is sparse and generically
invertible; in particular, at the exact specialization \((A,B,C)=(1,2,4)\) it
has full rank 30 and nonzero determinant.

Therefore the principal-symbol reduction has an explicit projector form. If
\[
\mathbf h^{(8)}\in\mathbb R^{45}
\]
collects the full degree-8 source coefficients in the ordered cartesian
monomial basis, then the selection matrices
\[
P_{\mathrm{orth}}\in\mathbb R^{15\times 45},
\qquad
P_{\mathrm{nonorth}}\in\mathbb R^{30\times 45}
\]
split the source as
\[
\mathbf h_{\mathrm{orth}}^{(8)}=P_{\mathrm{orth}}\mathbf h^{(8)},
\qquad
\mathbf h_{\mathrm{nonorth}}^{(8)}=P_{\mathrm{nonorth}}\mathbf h^{(8)}.
\]
Since \(i[S_{07},H_{02}]\) acts only on the nonorthorhombic complement, the
15 orthorhombic degree-8 coefficients are unchanged by this elimination step:
\[
\mathbf c_{\mathrm{orth}}^{(8)} = P_{\mathrm{orth}}\mathbf h^{(8)}
\]
at principal-symbol level. The role of \(S_{07}\) is only to solve
\[
A\mathbf s = -\,\mathbf h_{\mathrm{nonorth}}^{(8)}
\]
and remove the 30 nonorthorhombic components.

As a consequence, the direct principal-symbol source for the 15 orthorhombic
coefficients is already reduced from seven blocks to six:
\[
\bigl(
H08_{\mathrm{k4}},
S03S03S03H02,
S03S03H04,
S03H06,
S05H04,
S05S03H02
\bigr),
\]
while \(S07H02\) plays only the role of elimination block for the
nonorthorhombic complement.

The channel content is summarized in Table~\ref{tab:octic_channels}.

\begin{table}[t]
\caption{Principal-symbol channels entering the optical octic source.}
\label{tab:octic_channels}
\begin{adjustbox}{max width=\columnwidth}
\begin{tabular}{lllll}
\toprule
Block & Content & Primitive ingredients & Direct to \(L\) & Role \\
\midrule
\(H08_{\mathrm{k4}}\) & \(\Phi_4\,\mu_1^4\) & \(k'_{klmn},\,C_k^{ab}\) & yes & visible quartic source \\
\([S_{03},H_{06}]\) & \(\mu_1\,\Phi_3^2\) & \(\mu_1,\,\Phi_3,\,\omega,\,B\) & yes & direct cubic channel \\
\([S_{03},[S_{03},H_{04}]]\) & \(\mu_1\,\Phi_3^2\) & \(\mu_1,\,\Phi_3,\,\omega\) & yes & direct cubic channel \\
\([S_{05},H_{04}]\) & \(\mu_1\,\Phi_3\,\Phi_4\) & \(\mu_1,\,\Phi_3,\,\Phi_4,\,\omega\) & yes & mixed direct channel \\
\([S_{05},[S_{03},H_{02}]]\) & \(\Phi_3^2\,\Phi_4\) & \(\Phi_3,\,\Phi_4,\,\omega,\,B\) & yes & vibrational direct channel \\
\([S_{03},[S_{03},[S_{03},H_{02}]]]\) & \(\Phi_3^3\) & \(\Phi_3,\,\omega,\,B\) & yes & cubic direct channel \\
\([S_{07},H_{02}]\) & elimination only & \(\Phi_3,\,\Phi_4,\,\omega,\,B\) & no & removes nonorth complement \\
\bottomrule
\end{tabular}
\end{adjustbox}
\end{table}

The first channel that can already be developed into an explicit commuting
degree-8 source is the purely cubic block
\([S_{03},[S_{03},[S_{03},H_{02}]]]\). At principal-symbol level the
orthorhombic part of \(S_{03}\) is one-dimensional, with commuting symbol
\[
S_{03}^{(\mathrm{ps},\mathrm{orth})}
=
\sigma_3 J_xJ_yJ_z,
\qquad
\sigma_3 = 6 s_{xyz} = -\frac{3}{2}S_{111},
\]
where the last equality follows from the Watson normalization
\(s_{xyz}=-S_{111}/4\). Its triple Lie--Poisson action on
\(H_{02}=A J_x^2+B J_y^2+C J_z^2\) produces an explicit degree-8
orthorhombic polynomial, hence a concrete source term for the 15 octic
coefficients rather than a mere channel assignment.

The same mechanism applies to the quartic-cubic channel
\([S_{03},[S_{03},H_{04}]]\). Once the orthorhombic principal symbol of
\(H_{04}\) is written as the general six-parameter quartic commuting
polynomial,
\[
H_{04}^{(\mathrm{ps},\mathrm{orth})}
=
q_{400}J_x^4+q_{040}J_y^4+q_{004}J_z^4
+q_{220}J_x^2J_y^2+q_{202}J_x^2J_z^2+q_{022}J_y^2J_z^2,
\]
the double bracket with the one-dimensional orthorhombic \(S_{03}\) symbol
again yields an explicit degree-8 orthorhombic source. The
\(\mu_1\Phi_3^2\) block is therefore also concrete at principal-symbol level.

The sextic-cubic block \([S_{03},H_{06}]\) behaves in the same way.  Writing
\(H_{06}\) as the general 10-parameter orthorhombic sextic commuting
polynomial immediately yields, after one bracket with the one-dimensional
orthorhombic \(S_{03}\) symbol, an explicit degree-8 orthorhombic source.
Thus three of the six direct channels are already reduced concretely at
principal-symbol level. In this reduction the totally mixed sextic coefficient
\(w_{222}\) drops out identically, while the edge sextic coefficients remain.

The nested vibrational block \([S_{05},[S_{03},H_{02}]]\) is equally explicit
once the orthorhombic sector of \(S_{05}\) is written in its natural
three-parameter basis
\[
S_{05}^{(\mathrm{ps},\mathrm{orth})}
=
s_{311}J_x^3J_yJ_z+s_{131}J_xJ_y^3J_z+s_{113}J_xJ_yJ_z^3.
\]
Bracketing this against the explicit \([S_{03},H_{02}]\) source produces
another concrete degree-8 orthorhombic polynomial. Thus four of the six direct
channels are already explicit at principal-symbol level.

The last direct block \([S_{05},H_{04}]\) closes in exactly the same way:
the same three-parameter \(S_{05}\) source, bracketed with the general
orthorhombic quartic source \(H_{04}\), yields one more explicit degree-8
orthorhombic polynomial. At this point all six direct channels entering the
optical octic source are concrete at principal-symbol level.

The identification with the printed lower-order Watson--Aliev coefficients is
also already exact for the odd generators and for the useful sextic sector.
From Eq.~(89), \(\sigma_3=-3S_{111}/2\). From Eq.~(95), commuting reduction
gives \(s_{311}=-2S_{311}\), \(s_{131}=-2S_{131}\), and
\(s_{113}=-2S_{113}\). From Eq.~(96), the sextic principal coefficients obey
\(w_{600}=W_{xxx}\), \(w_{060}=W_{yyy}\), \(w_{006}=W_{zzz}\), and
\(w_{420}=2W_{xxy}\), \(w_{402}=2W_{xxz}\), \(w_{240}=2W_{yyx}\),
\(w_{042}=2W_{yyz}\), \(w_{204}=2W_{zzx}\), \(w_{024}=2W_{zzy}\). The only
remaining printed-to-principal link was the quartic \(H_{04}\) sector, and it
is equally simple at principal-symbol level:
\[
q_{400}=\frac14\tau_{xxxx},\qquad
q_{040}=\frac14\tau_{yyyy},\qquad
q_{004}=\frac14\tau_{zzzz},
\]
\[
q_{220}=\frac32\tau_{xxyy},\qquad
q_{202}=\frac32\tau_{xxzz},\qquad
q_{022}=\frac32\tau_{yyzz}.
\]
Thus the printed-to-principal map is closed for all lower-order objects needed
by the six direct octic channels.

At this point the six direct channels can be summed with the exact Eq.~(99)
prefactors to obtain the final orthorhombic degree-8 coefficients. The whole
\(S_{07}\) branch disappears from this final expression because it acts only on
the nonorthorhombic complement. The linear optical constant is the terminal
one-dimensional reduction of this orthorhombic octic sector, but its exact
operator-level relation to the raw cartesian coefficient \(c_{080}\) still has
to be established.
At the principal-symbol level developed so far, the five explicit commutator
channels have vanishing raw \(c_{080}\) component, so that
\[
c_{080}^{(\mathrm{ps})}=h08k4_{080}.
\]
From the visible quartic source term this coefficient is
\[
h08k4_{080}
=
\frac{1}{24}\sum_{klmn}
k'_{klmn}\,
C_k^{yy}C_l^{yy}C_m^{yy}C_n^{yy},
\]
since the \(J_y^8\) monomial is obtained by selecting the \(yy\) component of
each first-order rotational-derivative tensor.
This cancellation is not an artifact of the principal-symbol approximation:
an exact \(so(3)\) commutator audit with the full non-commuting rotational
algebra gives the same result for the raw \(J_y^8\) coefficient,
\[
c_{080}\bigl([S_{03},[S_{03},[S_{03},H_{02}]]]\bigr)
=
c_{080}\bigl([S_{03},[S_{03},H_{04}]]\bigr)
=
c_{080}\bigl([S_{03},H_{06}]\bigr)
=
c_{080}\bigl([S_{05},H_{04}]\bigr)
=
c_{080}\bigl([S_{05},[S_{03},H_{02}]]\bigr)
=0.
\]
This does not yet prove that the full spectroscopic constant \(L\) is exhausted
by the visible quartic source term, because the already established linear
formula for \(L\) also contains the \(-8D_J^3/B^2\), \(8BD_J C_n^2/\omega_n\),
and \(-r_n^2/(2\omega_n)\) pieces. The exact linear reduction is obtained only
after projecting the orthorhombic octic operator on the \(K=0\) sector and
reading the coefficient of \(X^4\), with \(X=J(J+1)\).

For the five ordered plane monomials of the linear \(a\)-axis convention,
\[
J_y^8,\qquad J_y^6J_z^2,\qquad J_y^4J_z^4,\qquad J_y^2J_z^6,\qquad J_z^8,
\]
the exact \(K=0\) reduction gives the \(X^4\) weights
\[
\frac{35}{128},\qquad \frac{5}{128},\qquad \frac{3}{128},\qquad \frac{5}{128},\qquad \frac{35}{128}.
\]
Therefore \(c_{080}=0\) does not imply a vanishing contribution to \(L\). The
commutator channels contribute nontrivially to the final \(X^4\) coefficient
after the exact linear projection. For example,
\[
[S_{03},[S_{03},[S_{03},H_{02}]]]
\;\longrightarrow\;
\frac{27}{64}S_{111}^3(C-B)\,X^4+\cdots,
\]
\[
[S_{03},[S_{03},H_{04}]]
\;\longrightarrow\;
\frac{9}{64}S_{111}^2(q_{004}-q_{022}+q_{040})\,X^4+\cdots.
\]
The same happens for the other three commutator channels. Thus the quartic
visible source alone does not exhaust the final spectroscopic \(L\), even
though it is the only direct source of the raw \(J_y^8\) coefficient.
The visible quartic source must itself be projected with the same \(K=0\)
weights. If
\[
h_{080},\ h_{062},\ h_{044},\ h_{026},\ h_{008}
\]
are its five plane coefficients in the linear \(a\)-axis convention, then its
contribution to the \(X^4\) coefficient is
\[
L^{(H08_{k4})}_{X^4}
=
\frac{35}{128}h_{080}
\,+\,
\frac{5}{128}h_{062}
\,+\,
\frac{3}{128}h_{044}
\,+\,
\frac{5}{128}h_{026}
\,+\,
\frac{35}{128}h_{008}.
\]
In particular, even if all mixed \(yz\) and \(zz\) pieces are artificially
suppressed, the surviving \(yy^4\) term contributes
\[
\frac{35}{128}\cdot \frac{1}{24}\sum_{klmn}k'_{klmn}
C_k^{yy}C_l^{yy}C_m^{yy}C_n^{yy},
\]
not the raw \(c_{080}\) coefficient itself.
After summing all six direct channels with the Eq.~(99) prefactors, the final
linear \(X^4\) coefficient is already nontrivial. Besides the projected
quartic \(H08_{k4}\) contribution, it contains explicit commutator pieces such
as
\[
\frac{9i}{128}(B-C)S_{111}^3,
\qquad
-\frac{9}{128}S_{111}^2(q_{004}-q_{022}+q_{040}),
\]
\[
\frac{3i}{128}S_{111}(-15w_{006}-w_{024}+w_{042}+15w_{060}),
\]
\[
\frac{i}{32}(5q_{004}s_{113}+3q_{004}s_{131}-q_{022}s_{113}+q_{022}s_{131}-3q_{040}s_{113}-5q_{040}s_{131}),
\]
\[
\frac{3}{32}(B-C)S_{111}(s_{113}-s_{131}).
\]
This is the first exact expression for the full linear \(X^4\) coefficient
obtained directly from the Watson--Aliev octic source hierarchy.
Moreover, in the true linear limit one must impose the transverse \(y/z\)
symmetry and quartic isotropy conditions
\[
C=B,\qquad q_{004}=q_{040},\qquad q_{022}=2q_{040},\qquad
w_{006}=w_{060},\qquad w_{024}=w_{042},\qquad s_{113}=s_{131},
\]
together with
\[
C_n^{zz}=C_n^{yy},\qquad C_n^{yz}=0.
\]
Under these exact linear-limit constraints, the whole principal-symbol
\(X^4\) coefficient collapses to
\[
\frac{1}{24}\sum_{klmn}k'_{klmn}C_k^{yy}C_l^{yy}C_m^{yy}C_n^{yy},
\]
which is precisely the quartic piece of the established linear formula for
\(L\). Therefore the additional
\[
-8D_J^3/B^2,\qquad 8BD_JC_n^2/\omega_n,\qquad -r_n^2/(2\omega_n)
\]
pieces must originate from the sub-principal terms of the one-line \(H_{08}\)
source, not from the principal-symbol channels treated above.
In the minimal commuting \(K=0\) matching ansatz, the one-line sub-principal
terms have exactly the required structure:
\[
RRR\,R_{klm}\ \longrightarrow\ C_kC_lC_m\,t_{klm}\,X^4,
\]
\[
R'_kRRR'_{lm}\ \longrightarrow\ -\frac{C_lC_m\,r_{lm}\,u_k}{\omega_k}\,X^4,
\]
\[
(R'_k)^2\ \longrightarrow\ -\frac{u_k^2}{2\omega_k}\,X^4.
\]
Thus the one-line \(H_{08}\) source has the exact algebraic capacity to produce
the missing \(D_J^3\), \(D_J C_n^2/\omega_n\), and \(r_n^2/(2\omega_n)\)
pieces of the linear formula for \(L\). The remaining task is to identify the
effective linear-limit coefficients \(t_{klm}\) and \(u_k\) from the Watson--
Aliev definitions of \(R_{klm}\) and \(R'_k\).
One point is already fixed without ambiguity:
\[
u_k=r_k.
\]
With this identification, the one-line contribution
\[
-\frac{(R'_k)^2}{2\omega_k}
\]
reduces exactly to
\[
-\frac{r_k^2}{2\omega_k},
\]
which is the corresponding term in the known linear formula for \(L\). The
only unmatched pieces are then
\[
8BD_J\frac{C_k^2}{\omega_k}
\qquad\text{and}\qquad
-8D_J^3/B^2,
\]
which must be carried by the remaining one-line terms
\[
RRR\,R_{klm}
\qquad\text{and}\qquad
R'_kRRR'_{lm}.
\]
The diagonal two-index prime block is also fixed. Using the screenshot-backed
relation
\[
R'_{kl}=R_{kl}-\frac12\sum_m \frac{k'_{klm}R_m}{\omega_m},
\]
its linear diagonal reduces to
\[
R'_{nn}=r_{nn},
\]
because the geometric part comes from \(R_{nn}\) and the cubic correction comes
from the explicit \(k'_{nnm}R_m/\omega_m\) subtraction. This is the second
exact bridge between the one-line \(H_{08}\) source and the already
established linear families.
However, the naive diagonal reduction of the one-line term
\[
R'_kRRR'_{lm}
\]
does \emph{not} reproduce the geometric piece
\[
8BD_J\frac{C_n^2}{\omega_n}.
\]
So the geometric contribution must come from a nontrivial effective reduction
of the full \(lm\) block, not from the diagonal \(R'_{nn}=r_{nn}\) alone.
If one writes the full linear reduction of that block in the effective form
\[
R'_kRRR'_{lm}\ \longrightarrow\ -\frac{r_k\,S_k[lm]}{\omega_k}\,X^4,
\]
then the exact matching condition is
\[
S_k[lm]
=
-\,8BD_J\,\frac{C_k^2}{r_k}.
\]
Under a minimal diagonal-only contraction ansatz this would reduce to
\[
\rho_k=-\,\frac{8BD_J}{r_k},
\]
but the previous mismatch shows that the true Watson--Aliev reduction cannot be
read off from the diagonal block alone. The remaining derivation problem is
therefore sharply localized: one must derive the full effective \(lm\)
contraction \(S_k[lm]\) from the printed one-line source.
For a single-mode probe this can be sharpened further. If \(u_n=r_n\) and
\(R'_{nn}=r_{nn}\), then the diagonal part of the contraction is
\[
S_n^{\mathrm{diag}}=C_n^2\,r_{nn},
\]
whereas the exact target is
\[
S_n^{\mathrm{target}}=-\,8BD_J\,\frac{C_n^2}{r_n}.
\]
Therefore the full Watson--Aliev \(lm\) contraction must contain the explicit
off-diagonal residue
\[
S_n^{\mathrm{off}}
=
S_n^{\mathrm{target}}-S_n^{\mathrm{diag}}
=
-\,8BD_J\,\frac{C_n^2}{r_n}-C_n^2r_{nn}.
\]
This is the precise piece that still has to be derived from the complete
non-diagonal reduction of \(R'_{lm}\).
Using the printed relation
\[
R'_{lm}=R_{lm}-\frac12\sum_n \frac{k'_{lmn}R_n}{\omega_n},
\]
the full contraction itself splits exactly as
\[
S_k[lm]=G_k[R_{lm}]+Q_k[k_3,R_n].
\]
After the \(K=0\) scalar reduction, the second term is fixed by the printed
subtraction:
\[
Q_k[k_3,R_n]
=
\frac12\sum_{lmn} C_lC_mC_n\,k'_{lmn}.
\]
So the residual derivation problem is now divided sharply into:
\begin{enumerate}
\item the known cubic-subtraction channel \(Q_k\);
\item the still unresolved geometric contraction \(G_k[R_{lm}]\), which must
carry the off-diagonal residue above.
\end{enumerate}
There is also a structural restriction. Since the printed one-line term
\[
R'_kR_lR_mR'_{lm}
\]
contracts only over \(l,m\) inside \(R'_{lm}\), the reduced scalar
contraction is independent of the external mode label:
\[
S_k[lm]\equiv S[lm].
\]
Therefore exact matching of
\[
8BD_J\frac{C_k^2}{\omega_k}
\]
through
\[
-\frac{r_k\,S[lm]}{\omega_k}
\]
requires
\[
S[lm]=-\,8BD_J\,\frac{C_k^2}{r_k}
\]
for every \(k\). Hence the ratio
\[
\frac{C_k^2}{r_k}
\]
must be mode-independent if this term alone is to reproduce the known linear
geometric contribution mode by mode.
But in the already established linear branch,
\[
r_k=4\omega_k C_k\!\left(\frac{D_J}{B}-\frac{B^2}{\omega_k^2}\right)+\text{cubic}_k,
\]
so
\[
\frac{C_k^2}{r_k}
\]
is generically not mode-independent. This remains true even if the cubic
correction is switched off. Therefore the one-line term
\[
R'_kR_lR_mR'_{lm}
\]
cannot by itself reproduce the geometric contribution
\[
8BD_J\frac{C_k^2}{\omega_k}
\]
mode by mode in the known linear limit.
The fourth visible one-line term also vanishes in the true linear limit. If
the exact rotational blocks \(R_k\) and \(R_l\) collapse to scalar multiples of
the same transverse operator \(J_y^2+J_z^2\), then
\[
[R_k,R_l]=0
\]
identically, and therefore
\[
-i\sum_{kl}\frac{R'_kR_l[R_k,R_l]}{\omega_k^2\omega_l}
\]
vanishes in the exact operator reduction. So this term cannot supply either of
the two missing residual pieces.
The same kind of obstruction appears on the \(RRR\,R_{klm}\) side. Using the
printed geometric definition of \(R_{klm}\) and imposing the true linear
\(y/z\) symmetry,
\[
R_{klm}^{(\mathrm{geom,lin})}
=
-\,\frac{\omega_k\omega_l\omega_m}{2B^2}\,C_kC_lC_m\,(J_y^3+J_z^3).
\]
If this is combined naively with the linear geometric forms of the three
\(R_k\) factors, the one-line term
\[
-\,\frac{R_kR_lR_mR_{klm}}{\omega_k\omega_l\omega_m}
\]
reduces only to an odd transverse object proportional to
\[
X^3(J_y^3+J_z^3),
\]
not to a pure scalar \(X^4\). So the constant piece
\[
-\,8D_J^3/B^2
\]
cannot come from the naive commuting geometric reduction of this term alone.
In fact the obstruction is stronger: the exact \(K=0\) diagonal matrix element
of \(X^3(J_y^3+J_z^3)\) vanishes identically for all tested \(J\). Therefore
the naive true-linear geometric reduction of \(RRR\,R_{klm}\) contributes
nothing at all to the linear scalar \(L\).
The corresponding one-mode source probe sharpens the picture. If one inserts
the true-linear one-mode source forms directly into
\[
-\,\frac{R_kR_lR_mR_{klm}}{\omega_k\omega_l\omega_m},
\]
the exact \(K=0\) diagonal vanishes identically for all tested \(J\). So, at
least in this probe, the first visible one-line source term does not generate
the missing constant contribution.
By contrast, the one-mode exact source probe for
\[
-\,\frac{R'_kR_lR_mR'_{lm}}{\omega_k\omega_l\omega_m}
\]
produces a genuine polynomial in \(X\) of degree five, with nonzero \(X^5\)
and \(X^4\) coefficients. This means that the third visible source term is
still a viable origin for the residual linear pieces, but only after the
source-level \(X^5\) content is reduced to the final effective octic scalar.
However, its diagonal one-mode scaling is already incompatible with the known
geometric term. After substituting the diagonal linear family \(r_{11}\), the
probe gives
\[
X^4 \sim \frac{\omega^3 C^7}{B^2}
\]
when the cubic correction is switched off, whereas the target geometric term
has the shape
\[
8BD_J\frac{C^2}{\omega}.
\]
Therefore the diagonal one-mode part of the third visible source term cannot
be the origin of the geometric residual. The only remaining possibility inside
this term is the genuinely multimode, non-diagonal contraction structure.
The first off-diagonal two-mode source probe confirms this. In that probe the
same term still produces a degree-5 polynomial in \(X\), and its \(X^4\)
coefficient scales as
\[
\frac{C_1C_2\,r'_{12}\,r'_1}{\omega_1},
\]
not as a simple quartic source monomial. So the only viable origin of the
geometric residual inside the printed source is now a genuinely nontrivial
multimode identification of the reduced one-index and two-index prime blocks.
The probe also imposes a sharp matching condition:
\[
r'_1r'_{12}=\frac{512}{5}BD_J\frac{C_1}{C_2},
\qquad
r'_2r'_{12}=\frac{512}{5}BD_J\frac{C_2}{C_1},
\]
hence
\[
\frac{r'_1}{r'_2}=\frac{C_1^2}{C_2^2}.
\]
But the already established linear branch gives
\[
\frac{r_1}{r_2}
=
\frac{C_1\omega_1\left(D_J/B-B^2/\omega_1^2\right)+\cdots}
{C_2\omega_2\left(D_J/B-B^2/\omega_2^2\right)+\cdots},
\]
which is generically different. Therefore even the first bimode probe is
incompatible with identifying the reduced one-index prime block with the known
linear \(r_k\) family.

Among these six direct blocks, \([S_{03},H_{06}]\) is the first nontrivial
cubic channel. Since \(S_{03}\) is purely vibrational and \(H_{06}\) already
carries the sextic \(\mu_1\)-\(\Phi_3\) content, this block belongs to the
\(\mu_1\Phi_3^2\) channel and contributes directly to the orthorhombic octic
coefficients, hence also to the non-linear optical constants and to the linear
limit \(L\).

The same holds for \([S_{03},[S_{03},H_{04}]]\). Since \(H_{04}\) carries the
quartic rotational \(\mu_1\) content and the two \(S_{03}\) generators supply
two cubic vibrational factors, this block also belongs to the
\(\mu_1\Phi_3^2\) channel and contributes directly to the optical constants and
to \(L\).

The mixed block \([S_{05},H_{04}]\) is likewise a direct degree-8 source term.
Since \(H_{04}\) carries the quartic rotational \(\mu_1\) content and
\(S_{05}\) is the quintic pure-vibrational generator built from the cubic and
quartic potentials, this block belongs to the \(\mu_1\Phi_3\Phi_4\) channel
and contributes directly to the optical constants and to \(L\).

Finally, the nested commutator \([S_{05},[S_{03},H_{02}]]\) is also a direct
degree-8 source term. Since \(H_{02}\) is purely harmonic rotational while
\(S_{03}\) and \(S_{05}\) are pure-vibrational generators, this block belongs
to the purely vibrational \(\Phi_3^2\Phi_4\) channel and contributes directly
to the optical constants and to \(L\) without introducing additional
rotational-derivative tensors.

The last direct block is the triple commutator
\([S_{03},[S_{03},[S_{03},H_{02}]]]\). Since \(H_{02}\) is purely harmonic
rotational and \(S_{03}\) is the cubic pure-vibrational generator, this block
belongs to the purely vibrational \(\Phi_3^3\) channel and contributes
directly to the optical constants and to \(L\), again without introducing
additional rotational-derivative tensors.

At principal-symbol level the source simplifies further. In the visible
one-line \(H_{08}\) source, only the quartic-potential term
\[
\sum_{klmn}\frac{R_kR_lR_mR_n\,k'_{klmn}}{24\,\omega_k\omega_l\omega_m\omega_n}
\]
already has degree 8, while the other four visible terms are subprincipal. By
contrast, all six completion blocks of Eq.~\eqref{eq:paper2_octic_eq99} are
degree-8 contributions at principal-symbol level. The degree-8 octic source
entering the orthorhombic projector is therefore assembled from that quartic
source term together with the six completion blocks:
\[
\mathbf h^{(8)}
=
\mathbf h^{(8)}_{\mathrm{H08\_k4}}
-\frac{i}{6}\mathbf h^{(8)}_{\mathrm{S03S03S03H02}}
-\frac12\mathbf h^{(8)}_{\mathrm{S03S03H04}}
+i\mathbf h^{(8)}_{\mathrm{S03H06}}
+i\mathbf h^{(8)}_{\mathrm{S05H04}}
-\mathbf h^{(8)}_{\mathrm{S05S03H02}}
+i\mathbf h^{(8)}_{\mathrm{S07H02}}.
\]
The first concrete contribution in this list is the visible quartic source
term. Using
\[
R_k=-\omega_k\sum_{ab} C_k^{ab}J_aJ_b,
\]
one obtains exactly
\[
\sum_{klmn}\frac{R_kR_lR_mR_n\,k'_{klmn}}{24\,\omega_k\omega_l\omega_m\omega_n}
=
\frac{1}{24}
\sum_{klmn}k'_{klmn}
\sum_{ab\,cd\,ef\,gh}
C_k^{ab}C_l^{cd}C_m^{ef}C_n^{gh}
J_aJ_bJ_cJ_dJ_eJ_fJ_gJ_h,
\]
so the harmonic frequencies cancel completely in this block. This is already a
pure \(\Phi_4\times\mu_1^4\) contribution to the non-linear optical source.
Among the remaining visible one-line source terms, the exact linear CeDiTT4
hierarchy now separates their roles sharply. Because the true-linear quartic
scalar space is one-dimensional and generated by
\((J_\perp^2)^2=J_y^4+2J_y^2J_z^2+J_z^4\), the visible contribution
\(- (R'_k)^2/(2\omega_k)\) is automatically a pure octic \(X_0\) term once
\(R'_k\) is recognized as a quartic transverse scalar. The visible
\(-iR'_kR_l[R_k,R_l]/(\omega_k^2\omega_l)\) term vanishes in the true linear
limit, and the naive \(R_kR_lR_mR_{klm}\) source has vanishing \(K=0\)
diagonal. The only visible subprincipal source term that survives as a
nontrivial hierarchy mixture is therefore
\(-R'_kR_lR_mR'_{lm}/(\omega_k\omega_l\omega_m)\).
At the rigorously isolated octic level, the visible source therefore
contributes to \(X_0\) through exactly three channels:
\begin{equation}
X_0^{\mathrm{vis}}
=
X_0^{(H08_{k4})}
+X_0^{(\mathrm{term}\;5)}
+X_0^{(\mathrm{term}\;3,\mathrm{red})},
\end{equation}
with term 1 and term 4 contributing zero in the true linear audit.
Combining this visible-source contribution with the completion-block total from
Eq.~\eqref{eq:paper2_octic_eq99} gives the exact derived octic scalar
\(X_0\) obtained so far in the present CeDiTT4 program.
At this stage the result can already be rewritten without the intermediate
principal-symbol coefficients \(q\), \(w\), \(s\), and \(\rho\): the total
derived \(X_0\) depends only on the primitive quartic tensor \(\tau\), the
printed sextic coefficients \(W\), the printed cubic generator coefficients
\(S\), the visible quartic \(H_{08}\) tensor block built from
\(k'_{klmn}C_k^{\alpha\beta}C_l^{\gamma\delta}C_m^{\epsilon\eta}C_n^{\mu\nu}\),
and the one-index linear scalar \(r_k\). In the true linear limit even this
last object can be reduced exactly from the printed definition of \(R'_k\):
since \([R_k,H_{02}]=0\) and the quartic commuting scalar space is
one-dimensional, one has
\[
r_k=\sum_l C_l r_{kl}.
\]
So the current total derived \(X_0\) may be written without \(r_k\) itself, in
terms of the contraction \(\sum_l C_l r_{kl}\). Using the printed reduction of
\(R'_{kl}\) once more, this contraction splits exactly into a geometric piece
\[
\frac{3}{4B}\,\omega_k\omega_l C_k C_l
\]
plus the scalar cubic correction coming from the
\(-\frac12\sum_m k'_{klm}R_m/\omega_m\) part of \(R'_{kl}\), namely
\[
\frac12\sum_m k'_{klm}C_m.
\]
So the current derived \(X_0\) is now written entirely in primitive tensorial
quantities.
If one also imposes the exact true-linear transverse constraints
\[
C_n^{zz}=C_n^{yy},\qquad C_n^{yz}=0,
\]
together with the axial equalities
\[
C=B,\qquad S_{131}=S_{113},\qquad \tau_{zzzz}=\tau_{yyyy},\qquad
\tau_{yyzz}=\frac{2}{3}\tau_{yyyy},\qquad W_{zzz}=W_{yyy},\qquad
W_{zzy}=W_{yyz},
\]
the currently derived scalar collapses exactly to
\begin{equation}
X_0^{\mathrm{derived,lin}}
=
\frac{1}{24}\,k'_{klmn} C_k^{yy} C_l^{yy} C_m^{yy} C_n^{yy}
+\frac{18}{35}S_{111}^2\tau_{yyyy}
-\frac{63}{256}K_{\mathrm{term3}}
-\frac{1}{2\omega_k}
\left[
\sum_l C_l\left(
\frac{3}{4B}\omega_k\omega_l C_k C_l
+\frac12\sum_m k'_{klm}C_m
\right)
\right]^2.
\end{equation}
This is the compact true-linear form derived so far from the present tensorial
program alone.
The contracted square can be split exactly into three pieces,
\[
G=\sum_l C_l\frac{3}{4B}\omega_k\omega_l C_k C_l,\qquad
Q=\sum_l C_l\frac12\sum_m k'_{klm}C_m,
\]
so that
\[
-\frac{1}{2\omega_k}(G+Q)^2
=
-\frac{G^2}{2\omega_k}
-\frac{GQ}{\omega_k}
-\frac{Q^2}{2\omega_k}.
\]
Hence the current derived true-linear scalar is naturally organized into:
the visible quartic block, the \(S_{111}^2\tau\) block, the constant
\(-\frac{63}{256}K_{\mathrm{term3}}\), and three exact contributions from the square of the printed
\(R'_k\) reduction: geometric, mixed geometric-cubic, and purely cubic.

The appendix will list explicitly the 15 degree-8 orthorhombic cartesian
monomials selected by \(P_{\mathrm{orth}}\), the complementary 30
nonorthorhombic monomials of degree 8 selected by \(P_{\mathrm{nonorth}}\), the
30 nonorthorhombic monomials of degree 7 entering \(S_{07}\), and the full
reduction chain
\[
\mathbf h^{(8)}
\longrightarrow
\mathbf c_{\mathrm{orth}}^{(8)}
\longrightarrow
\bigl(L_J^{(4)},L_{J^3K},L_{J^2K^2},L_{JK^3},L_K^{(4)}\bigr)
\longrightarrow
L.
\]

\section{Projection to Symmetric-Top and Linear Forms}

\section{Exact Octic Operator Hierarchy}

\subsection{Asymmetric-top commuting octic space}

Let \(\mathcal O_8\) denote the commuting pure rotational operator space of
total degree eight in \(J_x,J_y,J_z\). The asymmetric-top space has dimension
\begin{equation}
\dim \mathcal O_8^{\mathrm{asym}} = 15.
\end{equation}

\subsection{Symmetric-top reduction}

For a symmetric top with axis \(a\), the commuting octic space reduces to
\begin{equation}
\dim \mathcal O_8^{\mathrm{sym}} = 5,
\end{equation}
with basis
\begin{equation}
J^8,\qquad J^6J_a^2,\qquad J^4J_a^4,\qquad J^2J_a^6,\qquad J_a^8.
\end{equation}

Writing the corresponding coefficients as
\(
L_J^{(4)}, L_{J^3K}, L_{J^2K^2}, L_{JK^3}, L_K^{(4)}
\),
the exact projection from the cartesian commuting coefficients is
\begin{align}
L_J^{(4)} &= c_{080}, \\
L_{J^3K} &= -4c_{080}+c_{260}, \\
L_{J^2K^2} &= 6c_{080}-3c_{260}+c_{440}, \\
L_{JK^3} &= -4c_{080}+3c_{260}-2c_{440}+c_{620}, \\
L_K^{(4)} &= c_{080}-c_{260}+c_{440}-c_{620}+c_{800}.
\end{align}

\subsection{Linear limit}

In the linear limit only one commuting octic operator survives,
\begin{equation}
\dim \mathcal O_8^{\mathrm{lin}} = 1,
\end{equation}
namely
\begin{equation}
(J_\perp^2)^4=(J_y^2+J_z^2)^4.
\end{equation}
Therefore
the reduced linear octic sector is one-dimensional. The final spectroscopic
constant \(L\) must therefore be obtained from this one-dimensional operator,
but it should not yet be identified directly with the raw cartesian coefficient
\(c_{080}\).

\section{Linear Optical Constant as a Control Limit}

This section will contain:
\begin{enumerate}
\item the exact derivation of the linear optical branch from the printed
Watson--Aliev equations;
\item the comparison with the published linear compact formula for \(L\);
\item the benchmark relation \(L B^2/D^3 \approx -8\);
\item the comparison with available experimental optical constants.
\end{enumerate}

\section{Representation Changes and Pseudoinverse Transport}

The optical representation change is linear. If \(\mathbf c\) is the
15-component cartesian coefficient vector and \(\mathbf L_{\mathrm{sym}}\) the
5-component symmetric-top optical vector, then
\begin{equation}
\mathbf L_{\mathrm{sym}} = M_{5\times 15}\,\mathbf c.
\end{equation}
The corresponding lift is given by the Moore--Penrose pseudoinverse
\begin{equation}
M^{+} = M^{\mathrm T}(MM^{\mathrm T})^{-1}.
\end{equation}

For the linear image, the embedding vector and its left inverse will be written
explicitly in the final version.

\section{\texorpdfstring{$\mu_2$}{mu2}-Free Structure of the Final Optical Hamiltonian}

Under the same structural assumptions already established for the standard
sextic branch, the final optical Hamiltonian does not depend on \(\mu_2\). The
explicit dependency graph and the corresponding theorem statement will be
inserted here in the final version.

\section{General Non-Linear Optical Constants}

This section will contain the projection of the general cartesian octic source
onto:
\begin{enumerate}
\item the five symmetric-top optical constants,
\item the one-dimensional linear limit,
\item the asymmetric-top optical parameter set appropriate to the chosen
effective Hamiltonian.
\end{enumerate}

\section{Computational Strategy}

This section will summarize:
\begin{enumerate}
\item extraction of the rotational derivative tensors from quantum chemical
force fields;
\item assembly of the octic cartesian source coefficients;
\item projection/lift with linear maps and pseudoinverses;
\item validation against linear benchmarks and available experimental optical
constants.
\end{enumerate}

\section{Results}

\subsection{Linear molecules}

Planned benchmark systems:
\begin{itemize}
\item HCN,
\item HCCD,
\item \ce{C2H2},
\item OCS,
\item additional linear systems with reliable experimental optical constants.
\end{itemize}

\subsection{Symmetric-top and non-linear prototypes}

This section will collect the first non-linear applications once the full
cartesian source coefficients have been written explicitly in terms of
\(\mu_1\), \(\Phi_3\), and \(\Phi_4\).

\section{Discussion}

The final discussion will emphasize three points:
\begin{enumerate}
\item the linear optical constant \(L\) provides a fully controlled test case
because it is not contaminated by the unresolved compact \(U\)-channel problem
of the vibrational quartic branch;
\item optical constants remain in the linear tensor-transport regime;
\item \(H_{22}\) is the first sector that breaks this linearity in the broader
centrifugal-distortion hierarchy.
\end{enumerate}

\section{Conclusions}

The present paper establishes the structure of the optical centrifugal
distortion problem directly from the printed Watson--Aliev Hamiltonian and uses
the linear case as a control limit for the general non-linear extension.

\bibliography{paper2}

\end{document}
